\documentclass[11pt]{article}
\usepackage[T1]{fontenc}
\usepackage[a4paper,margin=1in]{geometry}
\usepackage{amsmath,amssymb,amsthm}
\usepackage{booktabs}
\usepackage{graphicx}
\usepackage[hidelinks]{hyperref}
\usepackage{microtype}

\newtheorem{theorem}{Theorem}
\newtheorem*{sourceproblem}{Source Open Problem 2.2}
\newtheorem*{remark}{Remark}

\setlength{\parindent}{0pt}
\setlength{\parskip}{0.52em}

\title{Schwarz Rearrangement Can Increase\
a Thresholded Nonlocal Energy}
\author{Run \texttt{fa\_banach\_001} --- lane 17}
\date{12 August 2026}

\begin{document}
\maketitle

\begin{center}
\fbox{\parbox{0.9\textwidth}{
\textbf{Status: counterexample, subject to human review.}
For $N=1$, $p=3/2$, and $\delta=1$, an explicit nonnegative four-level
step function satisfies $I_1(u^*)-I_1(u)>43/1000$.  Thus the universal
Schwarz-rearrangement inequality asked in Open Problem~2.2 of
arXiv:1704.08077 is false.
}}
\end{center}

\section{Source problem}

Nguyen and Squassina study, for $\delta>0$, the functional
\begin{equation}\label{eq:functional}
 I_\delta(v)=
 \iint_{\{|v(y)-v(x)|>\delta\}}
 \frac{\delta^p}{|x-y|^{N+p}}\,dx\,dy
\end{equation}
and prove that it can increase under a single polarization.  They then ask
on arXiv PDF page~5 of arXiv:1704.08077v2~\cite{source}:

\begin{sourceproblem}
For $N\geq1$, is $I_\delta(u^*)\leq I_\delta(u)$ for every measurable
$u:\mathbb R^N\to\mathbb R^+$ and every $\delta>0$?
\end{sourceproblem}

\begin{center}
\includegraphics[width=0.96\textwidth]{figures/open_problem_crop.png}
\end{center}

The exponent $p$ is part of the preceding definition and is suppressed in
the notation.  A single choice of $N,p,\delta,u$ with the reverse strict
inequality disproves the universal formulation.

\section{Idea}

The threshold in~\eqref{eq:functional} is strict.  We therefore use a step
function whose adjacent levels differ by exactly $\delta=1$.  The singular
kernel is never integrated across a touching pair of regions whose value gap
is active, so the energy remains finite.  The levels $3$ and $4$ form a
high-value core inside a long level-$2$ buffer.  Moving that core off center
does not change the distribution of values, but it lowers the sum of the
long-range interactions selected by the nonconvex threshold.  Schwarz
rearrangement moves the core back to the center and increases the energy.

\section{Counterexample}

\begin{theorem}\label{thm:main}
For $N=1$, $p=3/2$, and $\delta=1$, there is a nonnegative measurable
$u:\mathbb R\to\mathbb R$ such that
\[
 I_1(u^*)>I_1(u).
\]
More precisely, the difference is greater than $43/1000$.
\end{theorem}

\begin{proof}
Endpoint values below are immaterial.  Define
\[
u(x)=
\begin{cases}
4,&x\in[8,10],\\
3,&x\in[7,8]\cup[10,11],\\
2,&x\in[-14,7]\cup[11,14],\\
1,&x\in[-15,-14]\cup[14,15],\\
0,&\text{otherwise}.
\end{cases}
\]
Its superlevel measures are
\[
\begin{array}{c|cccc}
t\text{ range}&[0,1)&[1,2)&[2,3)&[3,4)\\ \hline
|\{u>t\}|&30&28&4&2.
\end{array}
\]
Consequently its one-dimensional Schwarz rearrangement is
\[
u^*(x)=
\begin{cases}
4,&x\in[-1,1],\\
3,&x\in[-2,-1]\cup[1,2],\\
2,&x\in[-14,-2]\cup[2,14],\\
1,&x\in[-15,-14]\cup[14,15],\\
0,&\text{otherwise}.
\end{cases}
\]

Only pairs of levels differing by at least two contribute to $I_1$.  Such
regions are separated by a positive distance, so both energies are finite.
For disjoint intervals $[a,b]$ and $[c,d]$ with $b<c$, direct integration
gives
\begin{align}\label{eq:pair}
 \int_a^b\int_c^d (y-x)^{-5/2}\,dy\,dx
 =\frac43\left(
 \frac1{\sqrt{c-b}}+\frac1{\sqrt{d-a}}
 -\frac1{\sqrt{c-a}}-\frac1{\sqrt{d-b}}
 \right).
\end{align}
If $d=\infty$, the two terms involving $d$ are simply omitted.  The left
exterior half-line is handled by reflection.  Summing twice~\eqref{eq:pair}
over the unordered pairs of regions whose values differ by more than one
yields
\begin{align*}
I_1(u)={}&\frac{112}{15}+\frac{2\sqrt2}{3}+\frac{2\sqrt6}{3}
+\frac{4\sqrt{26}}{39}+\frac{8\sqrt{21}}{63}
-\frac{8\sqrt3}{9}-\frac{8\sqrt7}{21}-\frac{8\sqrt{22}}{33}
-\frac{16\sqrt{29}}{87},\\
I_1(u^*)={}&\frac{28}{3}+\frac{16\sqrt{15}}{45}
+\frac{16\sqrt{17}}{51}-\frac{8\sqrt3}{9}
-\frac{32\sqrt{13}}{39}-\frac{16\sqrt{29}}{87}.
\end{align*}
Thus, after cancelling the common terms,
\begin{align}\label{eq:delta}
\Delta:=I_1(u^*)-I_1(u)
={}&\frac{28}{15}
-\frac{32\sqrt{13}}{39}-\frac{2\sqrt6}{3}-\frac{2\sqrt2}{3}
-\frac{8\sqrt{21}}{63}-\frac{4\sqrt{26}}{39} \notag\\
&+\frac{8\sqrt7}{21}+\frac{8\sqrt{22}}{33}
+\frac{16\sqrt{17}}{51}+\frac{16\sqrt{15}}{45}.
\end{align}

For a fully rational sign certificate, use the lower bounds
\[
\sqrt7>2.64575,\quad \sqrt{22}>4.690415,\quad
\sqrt{17}>4.123105,\quad \sqrt{15}>3.872983
\]
in the positive terms of~\eqref{eq:delta}, and the upper bounds
\[
\sqrt{13}<3.605552,\quad \sqrt6<2.449490,\quad
\sqrt2<1.414214,\quad \sqrt{21}<4.582576,\quad
\sqrt{26}<5.099020
\]
in its negative terms.  Each bound is verified by squaring.  Substitution
gives the exact rational estimate
\[
 \Delta>
 \frac{1376155019}{31906875000}>\frac{43}{1000}>0.
\]
This proves the reverse strict inequality.
\end{proof}

\section{Checks, scope, and novelty}

The construction is genuinely measurable and nonnegative.  Its jumps are
not a hidden source of infinite energy: every touching pair of constant
regions differs by at most one, whereas the defining condition is
$|u(x)-u(y)|>1$.  Formula~\eqref{eq:pair} therefore accounts for every
nonzero contribution.  The accompanying \texttt{verification.py} independently
enumerates all contributing pairs and verifies both the numerical difference
and the rational lower certificate.

Theorem~\ref{thm:main} answers Open Problem~2.2 negatively under its universal
reading, at the explicit exponent $p=3/2$.  It does not settle the paper's
separate Open Problem~4.1 concerning P\'olya--Szeg\H{o} for distributional
Riesz fractional gradients.  Nor does it classify which other fixed
exponents or higher dimensions satisfy or fail the inequality.

Before promotion, the four lightweight run indexes were searched by exact
arXiv id, title, and the functional/rearrangement keywords.  A bounded search
through 12 August 2026 used the exact title and combinations of
``$I_\delta(u^*)$'', ``Schwarz rearrangement'', ``threshold'', and
``nonlocal functional''.  It found the source and papers on different
rearrangement energies, but no later proof or counterexample for Open
Problem~2.2.  The journal landing page's cited-by list likewise exposed only
unrelated applications.  This is a bounded novelty check, not a claim of
exhaustive bibliographic priority.

\begin{thebibliography}{9}
\bibitem{source}
H.-M. Nguyen and M. Squassina,
\emph{Some remarks on rearrangement for nonlocal functionals},
Nonlinear Analysis \textbf{162} (2017), 1--12;
arXiv:1704.08077v2.
\end{thebibliography}

\end{document}
